% ECMLPKDD 2025 — Demo Track
% Springer LNCS format
\documentclass[runningheads]{llncs}

\usepackage[T1]{fontenc}
\usepackage{graphicx}
\usepackage{amsmath,amssymb}
\usepackage{hyperref}
\usepackage{url}
\usepackage{booktabs}
\usepackage{microtype}
\usepackage{xcolor}
\usepackage{algorithm}
\usepackage{algorithmic}

% ── Document ────────────────────────────────────────────────────────────────
\begin{document}

\title{TRAKNN Explorer: an Interactive Application for\\
Rare Meteorological Pattern Discovery}

% \titlerunning{TRAKNN Explorer}

\author{%
  % Author names — fill in
  Guillaume Coulaud\inst{1}
  % \and Author Two\inst{2}
}

\authorrunning{G. Coulaud}

\institute{
  Inria, France\\
  \email{guillaume.coulaud@inria.fr}
}

\maketitle

% ────────────────────────────────────────────────────────────────────────────
\begin{abstract}
We present \textsc{Traknn Explorer}, an interactive web application for
exploring rare and anomalous patterns in daily gridded climate data.
The application exposes a trajectory-based $k$-nearest-neighbour~(\textsc{knn})
scoring algorithm that assigns each day a scalar \emph{rarity score}
measuring how far its surrounding weather trajectory is from all
historical analogues.
The interface, built with Streamlit, provides four tightly integrated
analytical views: (i)~a scoring pipeline that runs the full
preprocessing and scoring workflow; (ii)~an anomaly explorer that ranks,
maps, and classifies detected events; (iii)~a cluster analysis module
that groups the top anomalies with $k$-means and displays composites; and
(iv)~an analogue search that retrieves the closest historical matches for
any user-specified query period.
The tool is aimed at climate scientists and machine-learning researchers
who need a transparent, parameter-rich, and reproducible entry point into
spatiotemporal anomaly detection.
\end{abstract}

\keywords{Spatiotemporal anomaly detection \and
          Climate extremes \and
          Trajectory similarity \and
          $k$-nearest neighbours \and
          Interactive exploration}

% ────────────────────────────────────────────────────────────────────────────
\section{Introduction}
\label{sec:intro}

Daily gridded reanalysis products such as ERA5 contain multi-decadal records
of atmospheric variables—mean sea-level pressure, temperature, geopotential
height—sampled on grids with hundreds of thousands of points.
Detecting which days carry genuinely unusual spatial configurations, rather than
ordinary seasonal or synoptic variability, is a prerequisite for attribution
studies, compound-event analyses, and the validation of climate model extremes.

Most existing approaches either rely on univariate indices (which discard
spatial structure), fit parametric models (which impose distributional
assumptions), or require labelled reference events.
We build on the non-parametric, label-free trajectory-kNN framework
proposed in prior work~\cite{placeholder}, in which a \emph{trajectory}
is defined as a sequence of $L$ consecutive spatial fields, and rarity is
quantified by the mean distance to the $k$ most distant non-adjacent
historical trajectories.

\textsc{Traknn Explorer} makes this framework accessible through an
interactive web interface.
Users can load any NetCDF variable, tune preprocessing and algorithm
parameters, run the scoring pipeline, and immediately inspect results—all
without writing a single line of code.
The application is open-source and runs locally, keeping data on the
researcher's own infrastructure.

% ────────────────────────────────────────────────────────────────────────────
\section{The TRAKNN Algorithm}
\label{sec:algo}

Let $\{x_t\}_{t=1}^{T}$ be a time series of spatial fields, where
$x_t \in \mathbb{R}^{H \times W}$ is the gridded observation on day $t$.
A \emph{trajectory} of length $L$ starting at day $i$ is the concatenation
$\mathbf{t}_i = (x_i, x_{i+1}, \ldots, x_{i+L-1})$.

\paragraph{Preprocessing.}
Before scoring, the pipeline applies the following optional steps in order:
\begin{enumerate}
  \item \emph{Leap-day removal.} February~29 entries are discarded so that
        every year has exactly 365 days, yielding a consistent seasonal index.
  \item \emph{Seasonal cycle removal.} A 365-day climatological mean
        $\bar{x}^{(\text{doy})}$ is estimated and subtracted from each frame,
        leaving anomaly fields.
  \item \emph{Cosine-latitude weighting.} Each spatial field is multiplied
        element-wise by $\sqrt{\cos(\varphi)}$, where $\varphi$ is the
        latitude of each grid point, correcting for the convergence of
        meridians at high latitudes.
  \item \emph{Pixel-wise standardisation.} Each grid point is standardised
        to unit temporal variance.
  \item \emph{PCA dimensionality reduction.} An optional Randomized SVD
        step retains the principal components that explain at least 95\%
        of the total variance, reducing the spatial dimension before distance
        computation.
\end{enumerate}

\paragraph{Trajectory distance via sliding-window recurrence.}
The squared Euclidean distance between trajectories $\mathbf{t}_i$ and
$\mathbf{t}_j$ decomposes as:
\begin{equation}
  d^2(\mathbf{t}_i, \mathbf{t}_j)
  = \sum_{\ell=0}^{L-1} \|x_{i+\ell} - x_{j+\ell}\|^2.
  \label{eq:traj-dist}
\end{equation}
Computing all $\binom{T}{2}$ pairwise trajectory distances naïvely costs
$O(T^2 L D)$ operations.
The algorithm exploits a sliding-window recurrence: the distance between
trajectories shifted by one day satisfies
$d^2(\mathbf{t}_{i+1}, \mathbf{t}_{j+1}) = d^2(\mathbf{t}_i, \mathbf{t}_j)
- \|x_i - x_j\|^2 + \|x_{i+L} - x_{j+L}\|^2$,
reducing the per-step update from $O(LD)$ to $O(D)$ after the first column
is computed.
Pairwise spatial distances are computed in blocked matrix operations that
exploit symmetry and keep memory bounded.

\paragraph{Rarity scoring.}
For each trajectory $\mathbf{t}_i$, trajectories within a temporal exclusion
window of $\pm E$ days are masked, and the rarity score is defined as:
\begin{equation}
  s_i = \frac{1}{k} \sum_{j \in \mathcal{N}_k(i)} d^2(\mathbf{t}_i, \mathbf{t}_j),
  \label{eq:score}
\end{equation}
where $\mathcal{N}_k(i)$ is the set of indices of the $k$ nearest neighbours
of $\mathbf{t}_i$ outside the exclusion window.
A large $s_i$ indicates that no close historical analogue exists, flagging
the period as genuinely rare.

\paragraph{Algorithm variants.}
Three scoring modes are available:
\texttt{base}~(standard exclusion applied at scoring time),
\texttt{exclusion}~(stricter rolling exclusion per query batch), and
\texttt{interval}~(interval-level scoring for sub-seasonal analysis).

\paragraph{Implementation.}
The scoring engine is implemented in PyTorch, enabling transparent
CPU/GPU execution via a single \texttt{device} parameter.
Batched blocked matrix multiplication keeps the peak memory footprint
independent of $T$, which is critical for multi-decadal ERA5 datasets.

% ────────────────────────────────────────────────────────────────────────────
\section{Application Description}
\label{sec:app}

\textsc{Traknn Explorer} is a Streamlit application structured around
a persistent sidebar for global parameters and four analytical tabs
(Figure~\ref{fig:app}).
All intermediate results are stored in the Streamlit session state,
so that later tabs can operate on the output of earlier ones without
re-running expensive computations.

\begin{figure}[t]
  \centering
  \includegraphics[width=\linewidth]{figures/app_screenshot.png}
  \caption{%
    The \textsc{Traknn Explorer} interface.
    The sidebar (left) controls data loading, preprocessing, and algorithm
    parameters.
    The four tabs (right) cover the scoring pipeline, anomaly exploration,
    cluster analysis, and analogue search.}
  \label{fig:app}
\end{figure}

\subsection{Sidebar: Global Parameters}

The sidebar collects all parameters shared across tabs:
\begin{itemize}
  \item \textbf{Data:} path to a NetCDF file and the variable name
        (e.g.\ \texttt{msl}, \texttt{t2m}, \texttt{z500}); optional
        spatial bounding box (latitude/longitude range) and temporal
        slice (start/end year).
  \item \textbf{Preprocessing:} toggle switches for leap-day removal,
        seasonal-cycle subtraction, cosine-latitude weighting, pixel-wise
        standardisation, and PCA reduction.
  \item \textbf{Algorithm:} trajectory length $L$, number of neighbours $k$,
        exclusion zone $E$, algorithm variant, and compute device.
\end{itemize}
A two-level caching strategy (file I/O in stage~1; preprocessing in stage~2)
ensures that changing algorithm parameters does not trigger a disk read, and
that changing the preprocessing flags does not require re-running the scoring.

\subsection{Tab 1 — Scoring Pipeline}

This tab is the entry point.
Clicking \emph{Run Scoring} loads and preprocesses the data, then invokes
the trajectory-kNN engine.
Users who have already run the pipeline offline can alternatively upload a
pre-computed \texttt{.npz} file (containing \texttt{scores} and
\texttt{times} arrays) produced by the command-line script
\texttt{case\_studies/score.py}.

Once scores are available, the tab displays:
(i)~a summary row of statistics (number of time steps, maximum, mean, minimum
score);
(ii)~an interactive histogram of the score distribution; and
(iii)~an interactive time series of the normalised scores.
Scores can be downloaded as a CSV file for further analysis.

\subsection{Tab 2 — Anomaly Explorer}

The Anomaly Explorer allows users to investigate the highest-scoring events
in detail.
Users select the number of top events $N$ and choose between ranking by
absolute score or by monthly $z$-score (which standardises scores within
each calendar month, highlighting events that are extreme relative to their
seasonal context).

The tab presents:
\begin{itemize}
  \item An interactive \textbf{time series} with the rarity score in the
        background and the top-$N$ events marked as red dots.
  \item A \textbf{ranked table} listing event date, score, and a rarity
        classification (\emph{common}, \emph{rare}, or \emph{extreme})
        based on the empirical CDF thresholds $F \geq 0.90$ and
        $F \geq 0.99$, respectively.
  \item A \textbf{spatial map} for the selected event, rendered with
        Cartopy in the PlateCarree projection.
        The displayed field is the anomaly with respect to the 365-day
        climatology (raw pressure anomaly in hPa), not the preprocessed
        scoring input.
\end{itemize}
Clicking a row in the table synchronises the map view to that event.
The table is downloadable as CSV.

\subsection{Tab 3 — Cluster Analysis}

The Cluster Analysis tab groups the top-$N$ anomalous events into coherent
meteorological regimes using $k$-means clustering.
A sweep over a user-defined range $[K_{\min}, K_{\max}]$ is run, reporting
the mean silhouette score and inertia for each $K$.
The best $K$ (maximising mean silhouette) is highlighted automatically.

The user can then inspect, for any selected $K$:
\begin{itemize}
  \item \textbf{Composite spatial maps} — mean anomaly field over all events
        in each cluster, rendered on a geographic projection with a shared
        diverging colormap and a common colorbar.
  \item \textbf{Cluster membership bar chart.}
  \item \textbf{Temporal distributions} — monthly and yearly histograms of
        events in each cluster, helping reveal seasonal preferences.
  \item \textbf{Per-event silhouette values} — a horizontal bar chart sorted
        by silhouette, exposing marginal members or potential outliers.
  \item \textbf{Date table} — sorted by rarity score, downloadable as CSV.
\end{itemize}
The clustering step has its own independent preprocessing controls
(including a detrending option not available in the scoring tab), so users
can compare clustering with and without long-term trends.

\subsection{Tab 4 — Analogue Search}

The Analogue Search tab answers a targeted scientific question: given a
specific weather episode (query period), which past episodes most closely
resembled it?
Users specify a query start date and duration (in days); the engine computes
the trajectory distance between the query and all historical windows,
applying an exclusion zone to avoid trivially close neighbours.

Results include:
\begin{itemize}
  \item \textbf{Rarity assessment} of the query period: normalised score,
        empirical percentile, and classification label.
  \item \textbf{Ranked analogue table} listing start dates and summed
        trajectory distances, downloadable as CSV.
  \item \textbf{Side-by-side spatial maps} comparing the mean anomaly field
        of the query period with the selected analogue period.
  \item \textbf{Temporal distributions} (monthly and yearly) of analogue
        occurrences.
  \item \textbf{Annotated time series} showing the rarity score time series
        with the query period highlighted and analogue positions marked.
\end{itemize}

% ────────────────────────────────────────────────────────────────────────────
\section{Demonstration Outline}
\label{sec:demo}

The demonstration will use ERA5 mean sea-level pressure (\texttt{msl}) over
the Euro-Atlantic domain (1950–2023) as the primary dataset.
The walkthrough will proceed as follows:

\begin{enumerate}
  \item \textbf{Scoring.} Configure the sidebar (trajectory length $L=5$,
        $k=10$, exclusion zone $E=5$, seasonal cycle removed, cosine-latitude
        weighting applied) and run the scoring pipeline.
        Point out the bimodal score distribution and the seasonal clustering
        of high-scoring events in winter.
  \item \textbf{Anomaly exploration.} Switch to monthly $z$-score ranking to
        surface summer extremes.
        Click on the 2003 European heat wave and on the 2010 Russian heat wave
        to display their pressure anomaly maps.
  \item \textbf{Cluster analysis.} Run the cluster sweep for $K \in [2, 10]$.
        Inspect the discovered circulation regimes and their seasonal
        distributions.
  \item \textbf{Analogue search.} Set the query to the onset of the 2003 heat
        wave.
        Show that the closest historical analogues are drawn from the same
        circulation family and discuss their temporal distribution.
\end{enumerate}

% ────────────────────────────────────────────────────────────────────────────
\section{Related Work}
\label{sec:related}

Quantile-based thresholding of scalar climate indices (e.g.~temperature
exceedance) is the most common operational approach to climate extremes but
ignores spatial coherence~\cite{ipcc_wg1_2021}.
Matrix-profile-based anomaly detection~\cite{yeh2016matrix} targets
univariate time series and does not directly extend to spatial fields.
Isolation Forest~\cite{liu2008isolation} and Local Outlier
Factor~\cite{breunig2000lof} operate on fixed-length feature vectors and
require an explicit vectorisation of the spatial field, losing temporal
ordering.
Deep-learning approaches such as autoencoders have been applied to
spatiotemporal data~\cite{munir2019deepant}, but require large labelled or
well-behaved training sets and are less interpretable.

TRAKNN is closest in spirit to the analogue method widely used in
meteorology~\cite{van2018statistical,yiou2014anawege}, but it generalises it
by scoring every trajectory in the record rather than retrieving analogues only
for a pre-selected query.
By expressing trajectories as sequences of spatial fields and computing
distances in the full grid space (or a PCA subspace thereof), the method
retains spatial structure while remaining non-parametric and label-free.

% ────────────────────────────────────────────────────────────────────────────
\section{Conclusion}
\label{sec:conclusion}

We have presented \textsc{Traknn Explorer}, an interactive Streamlit
application built on the trajectory-kNN rarity scoring framework.
The tool covers the full analytical workflow from raw NetCDF ingestion to
cluster-level interpretation and targeted analogue retrieval, making it
suitable both as a research instrument and as a teaching aid for machine
learning in climate science.
Future extensions include support for multi-variable trajectories,
probabilistic rarity calibration, and ensemble comparison modes.

\paragraph{Availability.}
Source code, installation instructions, and a demo dataset are available at
\url{https://github.com/placeholder/spatiotemporal-anomalies}.
The application is licensed under the CeCILL Free Software License.

% ── References ───────────────────────────────────────────────────────────────
\begin{thebibliography}{99}

\bibitem{placeholder}
Coulaud, G., et al.:
Trajectory-kNN rarity scoring for spatiotemporal anomaly detection.
Preprint (2025)

\bibitem{ipcc_wg1_2021}
IPCC:
Climate Change 2021: The Physical Science Basis. Contribution of Working
Group I to the Sixth Assessment Report.
Cambridge University Press (2021)

\bibitem{yeh2016matrix}
Yeh, C.-C.M., et al.:
Matrix profile I: All pairs similarity joins for time series: a unifying view
that includes motifs, discords and shapelets.
In: IEEE ICDM, pp.~1317--1322 (2016)

\bibitem{liu2008isolation}
Liu, F.T., Ting, K.M., Zhou, Z.-H.:
Isolation forest.
In: IEEE ICDM, pp.~413--422 (2008)

\bibitem{breunig2000lof}
Breunig, M.M., Kriegel, H.-P., Ng, R.T., Sander, J.:
\textsc{lof}: Identifying density-based local outliers.
In: ACM SIGMOD, pp.~93--104 (2000)

\bibitem{munir2019deepant}
Munir, M., Siddiqui, S.A., Dengel, A., Ahmed, S.:
DeepAnT: A deep learning approach for unsupervised anomaly detection in time
series.
IEEE Access \textbf{7}, 1991--2005 (2019)

\bibitem{van2018statistical}
Van den Dool, H.:
Empirical Methods in Short-Term Climate Prediction.
Oxford University Press (2007)

\bibitem{yiou2014anawege}
Yiou, P.:
AnaWEGE: a weather generator based on analogues of atmospheric circulation.
Geoscientific Model Development \textbf{7}(2), 531--543 (2014)

\end{thebibliography}

\end{document}
